% =====================================================================
%  Part IV -- The computational problem.
%  Source: tier0.md I.3 (1103-1147), I.4 (1149-1220), I.5 (1222-1256),
%          I.6 (1258-1313), I.7 (1315-1347), I.8 (1349-1423),
%          I.9 items 4-6 (1433-1438), plus II.8 items 4 and 7
%          (1816-1817, 1825-1826) relocated here.
% =====================================================================
\section{The computational problem}
\label{part:computational}

% ---------------------------------------------------------------------
\subsection{Why nuclear networks are the computational bottleneck}
\label{sec:bottleneck}

\begin{equation*}
\frac{\dd\mathbf{Y}}{\dd t} = \nu\,\mathbf{R}(\mathbf{Y};\,T,\rho) \qquad
n_{\text{species}} \text{ equations}, \; n_{\text{rxn}} \text{ terms}
\end{equation*}

\subsubsection{The three costs, quantified}
\label{sec:threecosts}

\textbf{1. Stiffness forces implicit integration.} The stiffness ratio
$S = \max|\mathrm{Re}\,\lambda|/\min|\mathrm{Re}\,\lambda|$ reaches
$10^{12}$--$10^{15}$ here (\S\ref{sec:stiffness}; ``$S > 10^{15}$ is not
uncommon in astrophysics'' --- \citealt{hix1999b}). An explicit method needs
${\sim}S$ steps. Implicit is not an optimization, it is a necessity.

\textbf{2. Each implicit step costs a Jacobian build plus a linear solve.}
Newton iteration on the backward-Euler / BDF residual requires forming
$J = \nu\,\partial R/\partial Y$ ($n \times n$) and factorizing $(I - hJ)$.
Dense LU is $\approx 2n^3/3$ FLOPs, i.e.\ $n^3/3$ multiply--add pairs
\citep{golub2013}; the table counts the latter:

\begin{center}
\begin{tabular}{@{}l r r r@{}}
\toprule
\textbf{Network} & \textbf{$n$} &
\textbf{Dense LU (multiply--adds/factorization)} & \textbf{Ratio} \\
\midrule
\msn{80}  & 80  & $1.7\times10^{5}$ & $1\times$ \\
\msn{151} & 151 & $1.1\times10^{6}$ & \textbf{$6.7\times$} \\
\msn{204} & 204 & $2.8\times10^{6}$ & $17\times$ \\
\bottomrule
\end{tabular}
\end{center}

\textbf{Why sparsity does not rescue you.} The Jacobian is sparse
\emph{structurally} --- each species couples only to reaction partners --- but
LU factorization \textbf{fills in}: eliminating one variable creates couplings
between all its neighbours. Nuclear networks are strongly connected through free
nucleons --- \code{neut} and \code{h1} each appear in ${\sim}41\%$ of
\msn{80}'s 607 reactions (either in 72\%), and \citet{hix1999b} describe exactly
this structure as a ``doubly bordered, band diagonal'' matrix whose dense border
carries the free $n$, $p$, $\alpha$ --- so those rows and columns are nearly
dense and the fill-in propagates. Reordering helps but does not restore the
sparsity. On top of that, each \emph{step} takes several Newton iterations, and
each iteration may re-factorize.

\textbf{3. Every zone, every step.} A stellar model has $10^{3}$--$10^{4}$ zones
(\citealt{hix1999b}: ``hundreds to thousands'' for 1-D models) and takes of
order $10^{4}$--$10^{6}$ timesteps to core collapse. The multiplication is what
makes the network a runtime majority-holder in late burning stages (within the
network step itself, the matrix solve alone consumes 90$+$\% of the time ---
\citealt{hix1999b}).

\textbf{The compromise no one is happy with:} use a network small enough to
afford, and accept that it gets the composition --- and therefore \Ye{} ---
wrong.

% ---------------------------------------------------------------------
\subsection{Softwired networks: the compromise MESA makes}
\label{sec:softwired}

MESA's response is the \textbf{softwired} network: a fixed species list, with
reaction links assembled at runtime from whichever rates connect those species
(as opposed to \emph{hardwired} fixed lists or \emph{approx} lumped networks ---
\S\ref{sec:mesabbq}).

An immediate methodological consequence: \textbf{the reaction count is a
measured quantity, not a documented one.} You have to ask MESA what it built.
That is the entire purpose of the Fortran probe
(\code{src/mesa\_probes/probe.f90}) and the reconciliation of S6, which
established 607 / 1518 reactions with \code{MESA\_ONLY = 0}.

\subsubsection{The actual shape of the trade-off}
\label{sec:networkshape}

Computed from the repository's own isotope lists:

\begin{figure}[H]
\centering
\begin{tikzpicture}[
  font=\footnotesize,
  x=0.95cm, y=0.62cm,
  sp/.style={font=\footnotesize\ttfamily, inner sep=0.6pt, anchor=base west},
  only151/.style={circle, draw=resultrule, line width=0.5pt, fill=white,
                  inner sep=0pt, minimum size=1.6mm},
  zlab/.style={font=\footnotesize, anchor=east},
  leg/.style={font=\scriptsize, anchor=west},
  note/.style={font=\scriptsize\itshape, anchor=west, text=warnrule}]

% mesa_80-only marker, set INSIDE the species label so it can never be
% misread as belonging to the next species along the row.
\def\m{\,\raisebox{0.15ex}{\textcolor{warnrule}{$\bullet$}}}
% a run of open circles: start slot #1, row #2, count #3
\def\opencircles#1#2#3{%
  \foreach \k in {1,...,#3}{\node[only151] at (#1+\k-1, #2) {};}}

% ---- axes
\draw[line width=0.7pt] (-0.4,-0.4) -- (-0.4,12.5);
\draw[-{Latex[length=2.2mm]}, line width=0.7pt] (-0.4,-0.4) -- (11.5,-0.4);
\node[font=\footnotesize, anchor=west] at (11.65,-0.4) {$N$};
\node[font=\footnotesize, anchor=south] at (-0.75,12.5) {$Z$};
\node[font=\scriptsize, anchor=north] at (5.0,-0.75)
  {(neutron-rich direction $\longrightarrow$)};

% ---- rows: Z label, named species, then open circles
\node[zlab] at (-0.55,12) {30};
\node[sp] at (-0.25,12) {zn60\m};

\node[zlab] at (-0.55,11) {29};
\node[sp] at (-0.25,11) {cu59\m};

\node[zlab] at (-0.55,10) {28};
\node[sp] at (-0.25,10) {ni56}; \node[sp] at (0.75,10) {ni57};
\node[sp] at (1.75,10) {ni58}; \node[sp] at (2.75,10) {ni59};
\opencircles{3.9}{10.08}{6}

\node[zlab] at (-0.55,9) {26};
\node[sp] at (-0.25,9) {fe52}; \node[sp] at (0.75,9) {fe53};
\node[sp] at (1.75,9) {fe54}; \node[sp] at (2.75,9) {\_\_};
\node[sp] at (3.75,9) {fe56};
\opencircles{4.9}{9.08}{6}

\node[zlab] at (-0.55,8) {24};
\node[sp] at (-0.25,8) {cr48}; \node[sp] at (0.75,8) {cr49};
\node[sp] at (1.75,8) {cr50};
\opencircles{2.9}{8.08}{5}

\node[zlab] at (-0.55,7) {22};
\node[sp] at (-0.25,7) {ti44}; \node[sp] at (0.75,7) {ti45};
\node[sp] at (1.75,7) {ti46};
\opencircles{2.9}{7.08}{5}

\node[zlab] at (-0.55,6) {21};
\node[sp] at (-0.25,6) {sc43};
\opencircles{0.9}{6.08}{6}
\node[note] at (7.1,6.05) {$\leftarrow$ \msn{80} has ONE scandium};

\node[zlab] at (-0.55,5) {20};
\node[sp] at (-0.25,5) {ca39\m}; \node[sp] at (0.75,5) {ca40};
\node[sp] at (1.75,5) {ca41}; \node[sp] at (2.75,5) {ca42};
\opencircles{3.9}{5.08}{7}

\node[zlab] at (-0.55,4) {18};
\node[sp] at (-0.25,4) {ar35\m}; \node[sp] at (0.75,4) {ar36};
\node[sp] at (1.75,4) {ar37}; \node[sp] at (2.75,4) {ar38};
\opencircles{3.9}{4.08}{3}

\node[zlab] at (-0.55,3) {14};
\node[sp] at (-0.25,3) {si27}; \node[sp] at (0.75,3) {si28};
\node[sp] at (1.75,3) {si29}; \node[sp] at (2.75,3) {si30};
\opencircles{3.9}{3.08}{3}

\node[zlab] at (-0.55,2) {10};
\node[sp] at (-0.25,2) {ne18\m}; \node[sp] at (0.75,2) {ne19};
\node[sp] at (1.75,2) {ne20}; \node[sp] at (2.75,2) {ne21};
\node[sp] at (3.75,2) {ne22};
\opencircles{4.9}{2.08}{1}

\node[zlab] at (-0.55,1) {8};
\node[sp] at (-0.25,1) {o14\m}; \node[sp] at (0.75,1) {o15};
\node[sp] at (1.75,1) {o16}; \node[sp] at (2.75,1) {o17};
\node[sp] at (3.75,1) {o18};
\opencircles{4.9}{1.08}{1}

% ---- legend, below the axis so it cannot collide with any row
\node[leg] at (-0.4,-1.75)
  {\raisebox{0.15ex}{\textcolor{warnrule}{$\bullet$}}\ \ $=$ in \msn{80} only
   (6 species)};
\node[only151] at (-0.28,-2.42) {};
\node[leg] at (-0.05,-2.5) {$=$ in \msn{151} only (77 species)};
\node[leg] at (-0.4,-3.25) {(unmarked) $=$ in both (74 species)};
\end{tikzpicture}
\caption{The two networks on the $(N,Z)$ plane, computed from the repository's
own isotope lists. \msn{151} is not a superset: it buys neutron-rich reach by
dropping the proton-rich upper-$Z$ corner.}
\label{fig:nzplane}
\end{figure}

\textbf{The two networks are not ``small'' and ``large'' --- they are
differently shaped.} This is the most important structural fact in Tier 0, and
it is invisible unless you diff the lists:

\begin{itemize}
\item \textbf{\msn{151} is not a superset of \msn{80}.} Six species ---
      \code{o14}, \code{ne18}, \code{ar35}, \code{ca39}, \code{cu59},
      \code{zn60} --- are in \msn{80} and \emph{absent} from \msn{151}.
\item \msn{80} reaches \textbf{$Z = 30$ (Zn)} but is narrow in neutron number.
\item \msn{151} caps at \textbf{$Z = 28$ (Ni)} but extends deep into
      neutron-rich territory: \code{fe55}, \code{fe57}--\code{fe61},
      \code{ni60}--\code{ni65}, \code{ca43}--\code{ca49},
      \code{sc44}--\code{sc49}, \code{ti47}--\code{ti51},
      \code{mn52}--\code{mn56}.
\end{itemize}

So the 77 extra species \msn{151} buys are almost entirely
\textbf{neutron-rich}, purchased partly by \emph{dropping} the proton-rich
upper-$Z$ corner. That is a deliberate physics choice: neutron-rich isotopes are
exactly the $\beta$-decay and electron-capture partners that govern \Ye{}.
\textbf{\msn{151} is not ``more of the same'' --- it is a network reshaped
toward getting the weak sector right.}

\subsubsection{The measured consequence}
\label{sec:measuredconsequence}

\msn{80} contains only 6 of 9 electron-capture controllers and \textbf{0 of 8}
$\beta$-decay partners of the \Ye{}-controller set. It is structurally incapable
of representing those \emph{dominant} neutron-rich $\beta$-decay chains ---
though not weak-one-sided overall: it still carries 19 $\beta^-$ channels and 19
\Ye{}-raising weak reactions in total (\code{RESULTS.md}, 2026-07-08).

Two things follow:

\begin{itemize}
\item The ``zero-shot \msn{80} $\to$ \msn{151} size transfer'' falsifier is not
      testing a superset extension. It tests whether a model trained on a
      network with \textbf{no} $\beta$-decay partners can predict a network
      dominated by them. Expect it to be hard, and know the reason.
\item The high-\Ye{} bottleneck $\nuc{Sc}{45}(p,\gamma)\nuc{Ti}{46}$ is a
      \textbf{\msn{151}-only} phenomenon. \msn{80} has exactly one scandium
      isotope, \nuc{Sc}{43}. You cannot even pose the question in \msn{80}.
\end{itemize}

\textbf{Predict the bias direction yourself} (self-check item~4 below): with EC
present but the dominant $\beta$-decay partners absent, \msn{80}'s \Ye{}
evolution is biased downward --- the restoring flux that remains is carried by
weaker channels.

% ---------------------------------------------------------------------
\subsection{Why machine-learning emulation}
\label{sec:whyml}

The network step is a \textbf{map}, and a well-posed one:

\begin{equation*}
\Phi_{\Delta t} : (T, \rho, \mathbf{X}) \;\mapsto\;
(\mathbf{X}', e_{\text{nuc}}, \varepsilon_\nu)
\end{equation*}

For fixed $(T, \rho)$ it is the time-$\Delta t$ flow of an \emph{autonomous} ODE
--- deterministic, smooth almost everywhere, fixed finite-dimensional input and
output. The autonomy is not free: it is granted by operator splitting, examined
properly in \S\ref{sec:splitting}.

\subsubsection{The cost model --- what speedup would actually matter}
\label{sec:costmodel}

Replace an adaptive implicit solve (several Newton iterations $\times$ several
substeps $\times$ $O(n^3)$ each) with one forward pass. For a GNN with $K = 5$
rounds and hidden width $d$ over $n$ nodes and $m$ edges, forward cost is
$O(K(n + m)d^2)$ --- linear in graph size, and \textbf{independent of
stiffness}.

That last clause is the real prize. Classical cost grows as the problem
stiffens; surrogate cost does not. In the hardest states --- precisely where the
network dominates runtime --- the ratio is most favourable.

\textbf{The strategic payoff:} if the surrogate is accurate, you can afford the
\emph{large} network everywhere, ending the accuracy-vs-cost compromise of
\S\ref{sec:softwired} rather than merely optimizing within it.

\subsubsection{Framing it honestly}
\label{sec:framinghonestly}

This is operator learning with a twist. The operator has exact algebraic
invariants (mass, charge, lepton number), and those invariants are not
decoration --- they are the observable (\S\ref{sec:explodability}). Note also
what is \emph{not} the difficulty: expressivity. An MLP can represent this map.
The difficulties are dynamic range, stiffness, exact constraints, and rollout
drift (\S\ref{sec:learningproblem}).

% ---------------------------------------------------------------------
\subsection{What breaks when an emulator does not conserve}
\label{sec:conservationfailure}

\subsubsection{Mode 1 --- accumulation}
\label{sec:accumulation}

Let the per-step conservation violation be $\delta$. Linearizing error
propagation about the true trajectory with contraction factor
$c = \|\partial\Phi/\partial X\|$ gives three regimes (derived in
\S\ref{sec:rollouterror}):

\begin{center}
\begin{tabular}{@{}l l c@{}}
\toprule
\textbf{Regime} & \textbf{Accumulated after $N$ steps} &
\textbf{log--log slope} \\
\midrule
Systematic ($c = 1$, correlated $\delta$)   & $N\delta$ & 1 \\
Random walk ($c = 1$, independent $\delta$) & $\sqrt{N}\,\delta$ & 0.5 \\
Contracting ($c < 1$) & $\delta/(1-c)$ --- saturates & $\to 0$ \\
\bottomrule
\end{tabular}
\end{center}

At $N \approx 10^{3}$--$10^{6}$ steps per stellar model, even
$\delta = 10^{-6}$ is fatal in the systematic case. This is exactly why the
project's per-step \Ye{} budget ($\lesssim 3\times10^{-6}$) is stated \emph{under
a stated accumulation model}, and why measuring the slope is the biggest single
lever.

\textbf{The prediction worth carrying} (\S\ref{sec:rollouterror}): silicon
burning contracts toward quasi-equilibrium, so \emph{composition} error should
saturate --- but \Ye{} has no restoring force, because strong reactions conserve
it exactly. Expect the two to behave differently, and expect the gate to belong
on \Ye{}.

\subsubsection{Mode 2 --- physical inconsistency}
\label{sec:physicalinconsistency}

$\sum_i X_i \neq 1$ --- equivalently $\sum_i A_i Y_i \neq 1$, equation
\eqref{eq:baryon-conservation} --- means the equation of state, the opacity, and
the energy generation are evaluated on a composition that does not exist. The
host code does not error; it silently produces a wrong star. (Note it is
$\sum X_i$ that is normalized, not $\sum A_iX_i$: the latter is
$\approx \bar{A} \sim 30$ and is not conserved by anything.)

\subsubsection{Why soft penalties are the wrong fix}
\label{sec:softpenalties}

The standard ML response is a loss term $\lambda\|C\cdot \dd Y\|^2$. This gives
conservation \emph{on average, on the training distribution, to within
optimization tolerance}. Every qualifier fails where you need it: rollouts leave
the training distribution, and averages do not constrain individual steps.
Worse, $\lambda$ becomes a hyperparameter trading accuracy against conservation
--- you are tuning how much physics to violate.

\textbf{Conservation by construction} removes the trade-off. If the decoder is
$\dd Y = \nu\varphi$ and $\nu$ is the exact integer stoichiometric matrix, then
$A\cdot\nu = 0$ \emph{identically}, so $A\cdot \dd Y = 0$ for \textbf{any}
$\varphi$ --- including a randomly initialized, untrained network. Conservation
becomes a property of the architecture, not of the training. (The same
soft-vs-hard contrast has been measured directly in climate-model emulation:
soft-constrained networks show a monotonic accuracy--conservation trade-off in
the penalty weight and large per-sample violations, while the architecturally
constrained variant conserves to machine precision without losing accuracy ---
\citealt{beucler2021}.)

This is why the conservation test asserts the gate holds \textbf{with random
flux heads}, and why it is a blocking gate rather than a metric. Note the
further precision this buys: floating-point analysis (\S\ref{sec:floatingpoint})
shows residual drift is bounded by ${\sim}n\varepsilon \approx
3.4\times10^{-14}$ in float64 --- 1.5 decades under the $10^{-12}$ gate. In
float32 it would miss by seven orders of magnitude.

% ---------------------------------------------------------------------
\subsection{Why a graph neural network}
\label{sec:whygnn}

The reaction network \emph{is} a graph --- not an analogy:

\begin{itemize}
\item \textbf{Nodes}: species (80 or 151), with features ($Z$, $A$, binding
      energy, abundance).
\item \textbf{Edges}: reactions; naturally \textbf{bipartite}, with species
      nodes and reaction nodes, since one reaction touches several species (the
      same representation NuGNN adopts --- \citealt{kim2026}).
\end{itemize}

Three properties follow --- formalism in \S\ref{sec:messagepassing}:

\begin{enumerate}
\item \textbf{Locality.} A reaction couples 2--4 species; information propagates
      in hops. Measured bipartite radius 3, diameter 6 --- and since one
      \emph{round} is species $\to$ reaction $\to$ species, i.e.\ \textbf{2
      bipartite hops}, $K = 3$ rounds already spans the graph. $K = 5$ is the
      project's choice and carries margin (\S\ref{sec:messagepassing}).
\item \textbf{Permutation equivariance.} The physics does not depend on species
      ordering; a GNN has this built in, an MLP must learn it, spending capacity
      on a symmetry that could have been free.
\item \textbf{Transferability.} Weights attach to node/edge \emph{types}, not to
      slots in a fixed-length vector --- which is what makes the
      \msn{80} $\to$ \msn{151} question askable at all, though
      \S\ref{sec:softwired} says why it is harder than it looks.
\end{enumerate}

\textbf{One design point worth noting now:} on the reaction $\to$ species
half-step, the correct aggregator is \textbf{sum}, because the true update is
literally a signed sum, $\dd Y_i = \sum_r \nu_{ir} \varphi_r$. Mean or max would
destroy extensivity. The architecture's aggregator choice is fixed by physics,
not by ablation.

The competing published surrogate (a dense feed-forward network) has none of
these; it is a fixed-width map from an $(80+2)$-vector to an $(80+2)$-vector ---
$(T, \rho, \mathbf{X})$ in, $(\mathbf{X}', e_{\text{nuc}}, \varepsilon_\nu)$ out
--- retrained from scratch for 151, with a separate network per timestep and
even a different depth per network size \citep{grichener2025}.

% ---------------------------------------------------------------------
\subsection{Where conservation is allowed to live}
\label{sec:whereconservation}

\subsubsection{Target A}
\label{sec:targetA}

Predict signed net per-reaction fluxes $\varphi$, decode through fixed
stoichiometry:

\begin{equation*}
\dd\mathbf{Y} = \nu\,\boldsymbol{\varphi}
\end{equation*}

Conserves for any $\varphi$. Also carries per-reaction structure: you can
inspect which reaction the model thinks is carrying flux, which is what makes
bridge and bottleneck analysis possible at all.

\subsubsection{Target B --- and the full projector derivation}
\label{sec:targetB}

Predict $\dd Y$ directly, then project onto the constraint manifold. We want the
\textbf{orthogonal projector onto $\ker C$}, where $C$ is the $m \times (n+3)$
constraint matrix (the construction is standard linear algebra --- e.g.\
\citealt{golub2013}).

Seek $P$ such that (i) $CP = 0$, (ii) $P v = v$ for $v \in \ker C$, (iii)
$P = P\Tr$. Decompose any vector into row-space and null-space parts. The row
space of $C$ is spanned by $C\Tr$, so write the row-space component as
$C\Tr\mathbf{a}$. Requiring $C(\mathbf{v} - C\Tr\mathbf{a}) = 0$ gives

\begin{equation*}
C\mathbf{v} = CC\Tr\mathbf{a} \quad\implies\quad
\mathbf{a} = (CC\Tr)^{-1} C\mathbf{v} \qquad
(CC\Tr \text{ invertible} \iff C \text{ full row rank})
\end{equation*}

Hence the null-space component is

\begin{equation}
\mathbf{v} - C\Tr(CC\Tr)^{-1}C\mathbf{v} \quad\implies\quad
\boxed{\;P = I - C\Tr(CC\Tr)^{-1}C\;}
\label{eq:projector}
\end{equation}

Verify the three properties:

\begin{itemize}
\item $\boldsymbol{CP} = C - CC\Tr(CC\Tr)^{-1}C = C - C = 0$ \yes
\item $\boldsymbol{P^2} = I - 2C\Tr(CC\Tr)^{-1}C
       + C\Tr(CC\Tr)^{-1}CC\Tr(CC\Tr)^{-1}C
       = I - C\Tr(CC\Tr)^{-1}C = P$ \yes
\item $\boldsymbol{P\Tr} = P$, since $(CC\Tr)^{-1}$ is symmetric \yes
\item $\operatorname{\mathbf{rank}} P = (n+3) - \operatorname{rank} C$ \yes
\end{itemize}

Note the requirement \textbf{$C$ must be full row rank} --- otherwise $CC\Tr$ is
singular. This is why \code{tests/test\_projector.py} includes a
rank-deficient-$C$ rejection test: it guards the existence condition of the
derivation, not just input validation. In practice the implementation uses a QR
null-space form, which is numerically better conditioned than forming
$(CC\Tr)^{-1}$ explicitly.

\subsubsection{The critical theorem --- and the documented failure mode}
\label{sec:criticaltheorem}

$P$ is valid \textbf{only in a linear output space}. Suppose you predict in a
signed-log space, $u = \asinh(\dd Y/s)$, and apply $P$ there. Then you have
enforced

\begin{equation*}
\sum_i A_i u_i = 0
\end{equation*}

which says \emph{nothing whatever} about $\sum_i A_i \dd Y_i$, because a sum
constraint does not commute with a nonlinear map. You get a beautifully
conserved quantity in a space with no physical meaning, and unconserved physics.

\textbf{Make the failure concrete.} Take two species with $A = 1$ and
$\dd Y = \pm10^{-3}$, so $\sum A\,\dd Y = 0$ exactly. Apply
$u = \asinh(\dd Y/s)$ with $s = 10^{-6}$: $u = \pm\asinh(10^{3}) \approx
\pm7.6$, and $\sum A u = 0$ too --- the symmetric case survives. Now take
$(+2\times10^{-3}, -10^{-3}, -10^{-3})$: $\sum A\,\dd Y = 0$ still, but
$\sum A u = 8.29 - 7.60 - 7.60 = -6.9 \neq 0$. Enforcing $\sum A u = 0$ would
therefore \emph{require} $\sum A\,\dd Y \neq 0$. The constraint sets are
genuinely different manifolds.

The rule that follows:

\begin{invariant}
\textbf{Conservation lives in the decode step, never inside latent dynamics.}
No nonlinear transform may sit between the conservation map and the output.
\end{invariant}

Internal $\asinh$/signed-log latents are fine --- they just may never be the
space where a sum constraint is evaluated.

% ---------------------------------------------------------------------
\subsection*{Self-check --- the computational problem}
\addcontentsline{toc}{subsection}{Self-check --- the computational problem}
\label{sec:selfcheck-computational}

\begin{selfcheck}
\item Why does sparsity fail to make the Jacobian factorization cheap? Name the
      specific species responsible.
\item Give the one-sentence argument for why soft conservation penalties cannot
      substitute for architectural conservation, referencing rollout.
\item Derive $P = I - C\Tr(CC\Tr)^{-1}C$ and verify idempotence. What breaks if
      $C$ is rank deficient, and which test guards it?
\item Using the diff in \S\ref{sec:networkshape}: list three $\beta$-decay
      chains representable in \msn{151} but not \msn{80}. Predict which
      direction \msn{80}'s \Ye{} evolution is biased.
      \emph{(Relocated here from the source document's \S II.8 item~4, which
      required \S I.4 --- the material now in \S\ref{sec:softwired}.)}
      \label{sc:mesadiff}
\item Construct a three-species example showing that $\sum A u = 0$ and
      $\sum A\,\dd Y = 0$ are different constraints under $u = \asinh(\dd Y/s)$.
      (Sketched in \S\ref{sec:criticaltheorem}.)
      \emph{(Relocated here from the source document's \S II.8 item~7, for the
      same reason.)}
\end{selfcheck}
